\documentclass[11pt]{article}
\usepackage[margin=1in]{geometry}
\usepackage{amsmath,amssymb,amsthm}
\usepackage[colorlinks=true,linkcolor=blue,urlcolor=blue,citecolor=blue]{hyperref}
\newtheorem{theorem}{Theorem}[section]
\newtheorem{lemma}[theorem]{Lemma}
\newtheorem{corollary}[theorem]{Corollary}
\newtheorem{remark}[theorem]{Remark}
\newcommand{\N}{\mathcal N}
\newcommand{\rank}{\operatorname{rank}}
\newcommand{\tr}{\operatorname{tr}}
\newcommand{\R}{\mathbb R}
\title{The sharp uniform coefficient for a doubled edge\protect\\
on finite unit-weight paths and cycles}
\author{Uniform path coefficient session}
\date{November 15, 2026}
\begin{document}
\maketitle

\begin{abstract}
For every fixed $0<\alpha<1$, the supremal strict eigenvalue count of
the energy-relative fractional-Laplacian correction for doubling one
edge on a finite unit-weight path, or on a connected graph of maximum
degree two, is asymptotic to
$\log^2(1/\varepsilon)/(2\pi^2)$.
The upper bound is uniform in the graph size and the edge location,
including arbitrary coupled growth with the inverse tolerance.
Its proof uses a finite-measure approximation lemma: periodization on
each logarithmic spectral block has the sharp Fourier rank, and the
periodization images have a separate small-rank exponential expansion.
Only the leverage-measure bound
$\nu((0,\delta])\le\sqrt\delta$ is needed from the finite graphs.
The centered-path limit supplies the lower bound.
\end{abstract}

\section{Statement}
All graphs below are finite, connected, and simple, with unit vertex
masses and unit edge conductances before the update. For an existing
edge $e=\{u,v\}$, put $b=e_u-e_v$, $B=L+bb^*$, and
\[
 C_\alpha=(L^\dagger)^{\alpha/2}
                  (B^\alpha-L^\alpha)(L^\dagger)^{\alpha/2},
 \qquad
 \N_C(\varepsilon)=\#\{j:\lambda_j(C)>\varepsilon\}.
\]
Pseudoinverse powers vanish on constants. Let
$W_{\mathrm{path},\alpha}$ be the supremum of this count over all paths
and edited edges, and let $W_{2,\alpha}$ allow all connected hosts of
maximum degree at most two. Degree refers to the combinatorial graph
before doubling the edge.

\begin{theorem}\label{thm:main}
For each fixed $0<\alpha<1$, as $\varepsilon\downarrow0$,
\[
 W_{\mathrm{path},\alpha}(\varepsilon)
 \sim W_{2,\alpha}(\varepsilon)
 \sim \frac{1}{2\pi^2}\log^2(1/\varepsilon).
\]
In fact, uniformly over every finite path and cycle and every edited
edge, with $L_\varepsilon=\log(1/\varepsilon)$,
\begin{equation}\label{eq:quantitative}
 \N_C(\varepsilon)
 \le \frac{L_\varepsilon^2}{2\pi^2}
                  +O_\alpha(L_\varepsilon^{3/2}).
\end{equation}
The constant in the remainder is independent of the number of
vertices and the edge location.
\end{theorem}

The previously established centered infinite-path coefficient alone
does not give the uniform upper bound. The uniform ingredient used
here is the finite-measure argument in
Sections~\ref{sec:block}--\ref{sec:tail}.
It does not require monotonicity of finite-path eigenvalues or any
limit comparison for an upper bound.

\section{The positive kernel and exact finite-graph measures}
Restrict first to the nonconstant subspace, and write $A=L\succ0$.
For $c_\alpha=\sin(\pi\alpha)/\pi$, the fractional resolvent formula
and the rank-one resolvent identity give
\begin{equation}\label{eq:resolvent}
 C_\alpha=c_\alpha\int_0^\infty
       \frac{t^\alpha v_tv_t^*}{1+m(t)}\,dt,
 \qquad v_t=A^{-\alpha/2}(A+tI)^{-1}b,
 \quad m(t)=b^*(A+tI)^{-1}b.
\end{equation}
Let $H_\alpha$ denote the same integral without the denominator.
If $\tau=b^*A^{-1}b$, then
\begin{equation}\label{eq:sandwich}
 (1+\tau)^{-1}H_\alpha\preceq C_\alpha\preceq H_\alpha.
\end{equation}
For an existing unit edge, $bb^*\preceq A$, so $\tau\le1$;
on a path, the edge is a bridge and $\tau=1$.

In an orthonormal eigenbasis $A\phi_k=\lambda_k\phi_k$, put
\[
 z_k=\frac{|b^*\phi_k|^2}{\lambda_k},\qquad
 \nu=\sum_k z_k\delta_{\lambda_k}.
\]
After diagonal phase conjugation and deleting zero weights, the matrix
of $H_\alpha$ is
\begin{equation}\label{eq:kernel}
 (H_\alpha)_{ij}=\sqrt{z_i z_j}
                 K_\alpha(\log\lambda_i-\log\lambda_j),
 \quad
 K_\alpha(s)=\frac{\sinh(\alpha s/2)}{\sinh(s/2)},
 \quad K_\alpha(0)=\alpha.
\end{equation}
This follows by evaluating the scalar divided-difference integral in
\eqref{eq:resolvent} without its denominator. Repeated eigenvalues
can either be retained as labelled atoms or combined: both descriptions
have the same nonzero operator spectrum. Thus $H_\alpha$ is the
integral operator with this kernel on $L^2(\nu)$.

The Fourier convention and symbol used throughout are
\begin{equation}\label{eq:fourier}
 \widehat K_\alpha(\xi)=\int_\R e^{-i\xi s}K_\alpha(s)\,ds
 =\frac{2\pi\sin(\pi\alpha)}
       {\cosh(2\pi\xi)+\cos(\pi\alpha)}.
\end{equation}
This is the rescaled $\sinh(at)/\sinh(\pi t)$ entry of NIST
DLMF Table~1.14.1~\cite{DLMF}. In particular $K_\alpha$ is positive
definite and, for a constant depending only on fixed $\alpha$,
\begin{equation}\label{eq:fourierdecay}
 0<\widehat K_\alpha(\xi)\le C_\alpha e^{-2\pi|\xi|}.
\end{equation}
Also $0<K_\alpha(s)\le\alpha$ for real $s$; the latter follows from
convexity of $\sinh$ on $[0,\infty)$.

\begin{lemma}[Finite path and cycle measures]\label{lem:measures}
For the path on $n$ vertices and edge $(p,p+1)$,
\[
 \lambda_k=4\sin^2\frac{k\pi}{2n},\qquad
 z_k=\frac2n\sin^2\frac{kp\pi}{n},\qquad 1\le k<n.
\]
For a cycle on $n\ge3$ vertices, a complex Fourier basis gives
\[
 \lambda_k=4\sin^2\frac{k\pi}{n},\qquad z_k=\frac1n,
 \qquad 1\le k<n.
\]
Consequently both families have spectrum in $(0,4]$ and satisfy
\begin{equation}\label{eq:cdf}
                 \nu((0,\delta])\le\sqrt\delta\qquad(\delta>0),
\end{equation}
uniformly in every graph size and edited edge.
\end{lemma}
\begin{proof}
On the path, subtract the values of the normalized eigenvector
$\phi_k(j)=\sqrt{2/n}\cos(k\pi(j-1/2)/n)$ at $p,p+1$, then divide
its squared modulus by $\lambda_k$. Since
$\sin x\ge2x/\pi$ on $[0,\pi/2]$, one has
$\lambda_k\ge4k^2/n^2$. Hence at most $n\sqrt\delta/2$ indices
can lie below $\delta$, each of weight at most $2/n$.
On a cycle, $|b^*\phi_k|^2=\lambda_k/n$ and
$\lambda_k\ge16\min(k,n-k)^2/n^2$. Counting both ends of the
Fourier grid even gives $\nu((0,\delta])\le\sqrt\delta/2$.
The same counting inequalities remain valid when $\delta\ge4$.
\end{proof}

\section{A uniform rank lemma on one logarithmic block}\label{sec:block}
\begin{lemma}[Periodization and its image correction]\label{lem:block}
Let $I\subset\R$ be an interval of length $h\ge1$. For every $R\ge0$
there is a real symmetric kernel $F_R(x,y)$ on $I\times I$ with
\begin{align}
 \sup_{x,y\in I}|K_\alpha(x-y)-F_R(x,y)|&\le e^{-R},
                                                        \label{eq:kernelerror}\\
 \rank F_R&\le
       \left(\frac{h+1}{2\pi^2}+4\right)R+C_\alpha(h+1).
                                                        \label{eq:blockrank}
\end{align}
Here rank is the number of separated terms needed for the kernel.
In particular, for every positive measure $\mu$ of mass $m$ supported
on $I$, its $K_\alpha(x-y)$ operator on $L^2(\mu)$ has an approximant
of this rank and operator-norm error at most $m e^{-R}$.
All constants are independent of $\mu$, including its number of atoms.
\end{lemma}
\begin{proof}
Put $T=h+1\ge2$ and periodize:
\[
 P_T(s)=\sum_{r\in\mathbb Z}K_\alpha(s+rT).
\]
Exponential decay of $K_\alpha$ makes this a continuous periodic
function. Integrating the periodization gives its Fourier coefficients
$T^{-1}\widehat K_\alpha(2\pi k/T)$. Their absolute summability and
\eqref{eq:fourierdecay} show that the truncation $P_{T,q}$ to
$|k|\le q$ has uniform error bounded by
\begin{equation}\label{eq:periodictail}
 \frac{2C_\alpha}{T}
 \frac{e^{-4\pi^2(q+1)/T}}{1-e^{-4\pi^2/T}}
 \le C'_\alpha e^{-4\pi^2q/T}.
\end{equation}
The prefactor is bounded uniformly for $T\ge2$: use
$1-e^{-v}\ge v/(1+v)$.
As a kernel in $s=x-y$, $P_{T,q}$ has real rank at most $2q+1$;
pair the positive and negative frequencies into sine and cosine terms.

It remains to remove the nonzero periodization images without losing
the leading Fourier rank. Set
\[
 a=\frac{1-\alpha}{2}>0,\qquad b=\frac{1+\alpha}{2}>a.
\]
For $t>0$ the geometric expansion gives exactly
\[
 K_\alpha(t)=\sum_{j=0}^\infty
                  \bigl(e^{-(j+a)t}-e^{-(j+b)t}\bigr).
\]
For $|s|\le h$ all nonzero image arguments stay at least one from
zero. Absolute convergence therefore permits summing first in the
image index, with result
\begin{equation}\label{eq:images}
 P_T(s)-K_\alpha(s)
 =\sum_{j=0}^\infty\left[
 d_{j+a}(e^{(j+a)s}+e^{-(j+a)s})
 -d_{j+b}(e^{(j+b)s}+e^{-(j+b)s})\right],
 \quad d_c=\frac{e^{-cT}}{1-e^{-cT}}.
\end{equation}
Keep $0\le j<J$ and call the resulting kernel $E_{T,J}$.
After substituting $s=x-y$, each retained $j$ contributes at most
four separated terms. It is a real symmetric kernel of rank at most
$4J$. Its tail satisfies
\begin{align*}
 \sup_{|s|\le h}|(P_T-K_\alpha)(s)-E_{T,J}(s)|
 &\le\sum_{j\ge J}\sum_{c\in\{j+a,j+b\}}
                  \frac{2e^{-c(T-h)}}{1-e^{-cT}}\\
 &\le C''_\alpha e^{-J},
\end{align*}
because $T-h=1$ and $T\ge2$.

Choose a fixed $c_\alpha\ge0$ large enough that
$(C'_\alpha+C''_\alpha)e^{-c_\alpha}\le1$, and set
\[
 q=\left\lceil\frac{T(R+c_\alpha)}{4\pi^2}\right\rceil,
 \qquad J=\lceil R+c_\alpha\rceil,
 \qquad F_R(x,y)=P_{T,q}(x-y)-E_{T,J}(x-y).
\]
The two tail estimates prove \eqref{eq:kernelerror}, and
$\rank F_R\le2q+1+4J$ proves \eqref{eq:blockrank}.
For an arbitrary measure of mass $m$, an error kernel bounded by
$e^{-R}$ has Hilbert--Schmidt norm at most $m e^{-R}$.
This implies the operator-norm assertion. This step is a norm
estimate; no pointwise kernel comparison is used to infer a
Loewner inequality.
\end{proof}

\section{A sharp upper coefficient from a measure tail}\label{sec:tail}
\begin{theorem}[Arbitrary measures with exponential tails]\label{thm:measure}
Fix $\alpha\in(0,1)$ and $A,\beta>0$. Let $\mu$ be any positive
measure on $[0,\infty)$ satisfying
\begin{equation}\label{eq:tail}
                      \mu([s,\infty))\le A e^{-\beta s}
                      \qquad(s\ge0).
\end{equation}
Let $T_\mu$ be the positive compact integral operator with kernel
$K_\alpha(x-y)$ on $L^2(\mu)$. Uniformly over all such measures,
\begin{equation}\label{eq:generaltailbound}
 \N_{T_\mu}(e^{-L})
 \le \frac{L^2}{4\pi^2\beta}+O_{\alpha,A,\beta}(L^{3/2})
 \qquad(L\longrightarrow\infty).
\end{equation}
\end{theorem}
\begin{proof}
The operator is compact since $\mu$ is finite and its kernel is
bounded; positivity follows from the positive Fourier symbol.
For a positive operator $T$ and $D$ complementary orthogonal block
projections $P_j$, Cauchy--Schwarz gives
\begin{equation}\label{eq:pinching}
 T\preceq D\sum_jP_jTP_j.
\end{equation}
Indeed $\|T^{1/2}\sum_jP_jf\|^2
\le D\sum_j\|T^{1/2}P_jf\|^2$.

Put $\varepsilon=e^{-L}$ and, for sufficiently large $L$, choose
\[
 h=\sqrt L,\qquad
 S=\frac{L+4\log(L+2)}\beta,\qquad
 M=\lceil S/h\rceil,\qquad D=M+1.
\]
Partition $[0,\infty)$ into the $M$ half-open intervals
$I_i=[ih,(i+1)h)$, $0\le i<M$, and the tail $[Mh,\infty)$.
Thus atoms on boundaries have a unique assigned block.
The mass $m_i$ of $I_i$ is at most $A e^{-\beta ih}$.
The tail compression has norm at most its mass times $\alpha$, hence
at most
\[
 \alpha A e^{-\beta Mh}
 \le \alpha A\varepsilon(L+2)^{-4}
 <\varepsilon/D
\]
for sufficiently large $L$, since $D=O_\beta(\sqrt L)$.
By \eqref{eq:pinching} and eigenvalue min--max,
\begin{equation}\label{eq:countblocks}
 \N_{T_\mu}(\varepsilon)
 \le\sum_{i=0}^{M-1}\N_{P_iT_\mu P_i}(\varepsilon/D).
\end{equation}

Set $U=L+\log(AD)$ and
\[
                         R_i=(U-\beta ih)_+.
\]
Lemma~\ref{lem:block} gives a block approximant of rank
\[
 r_i\le\left(\frac{h+1}{2\pi^2}+4\right)R_i+C_\alpha(h+1)
\]
with error at most $m_i e^{-R_i}\le\varepsilon/D$.
This inequality holds also when $R_i=0$, because then
$m_i\le A e^{-\beta ih}\le\varepsilon/D$.
For a positive compact operator, a rank-$r_i$ approximant at this
error implies strict count at most $r_i$. Therefore
\[
 \N_{T_\mu}(\varepsilon)
 \le\left(\frac{h+1}{2\pi^2}+4\right)\sum_{i=0}^{M-1}R_i
                                      +C_\alpha M(h+1).
\]
For large $L$, $U>0$. Comparison of the decreasing positive part
with its triangle integral yields
\[
 \sum_{i=0}^{M-1}R_i
 \le\sum_{i=0}^\infty(U-\beta hi)_+
 \le U+\int_0^\infty(U-\beta ht)_+\,dt
 =U+\frac{U^2}{2\beta h}.
\]
Now $U=L+O_{A,\beta}(\log L)$, $M=O_\beta(\sqrt L)$, and
$h=\sqrt L$. Substitution gives
\[
 \N_{T_\mu}(e^{-L})
 \le\frac{L^2}{4\pi^2\beta}
       +O_{\alpha,A,\beta}(L^{3/2}+L\log L),
\]
which is \eqref{eq:generaltailbound}.
Every estimate was uniform in the measure.
\end{proof}

\begin{remark}[The general tail coefficient is sharp]
For $d\mu(s)=e^{-\beta s}\,ds$ on $s\ge0$, passage to Lebesgue
measure gives the weighted convolution
$e^{-\beta(s+t)/2}K_\alpha(s-t)$. The Fourier symbols at $\alpha$
and $1/2$ are comparable above and below by fixed positive constants.
For $\alpha=1/2$, the unitary substitution $u=e^{-s/2}$ gives the
kernel $2u^\beta v^\beta/(u+v)$ on $(0,1)$.
Laptev's parameter is $\rho=2\beta$, so its count coefficient is
$1/(4\pi^2\beta)$. Thus the constant in
Theorem~\ref{thm:measure} cannot be reduced using just the tail
hypothesis, even for non-atomic measures. Multiplying the measure by
a fixed constant accommodates any prescribed positive tail prefactor
$A$ and does not change the leading coefficient.
\end{remark}

\begin{corollary}[General leverage-CDF coefficient]\label{cor:cdf}
Consider any family of positive rank-one updates with bounded positive
baseline spectra $\lambda\le\Lambda$ and with leverage measures
$\nu((0,\delta])\le K\delta^\beta$, uniformly in the family.
For fixed $0<\alpha<1$ its relative fractional correction satisfies
\[
 \N_C(e^{-L})\le\frac{L^2}{4\pi^2\beta}
                      +O_{\alpha,K,\Lambda,\beta}(L^{3/2}).
\]
\end{corollary}
\begin{proof}
Push $\nu$ forward under $s=\log(\Lambda/\lambda)$.
Its tail is at most $K\Lambda^\beta e^{-\beta s}$.
Equation~\eqref{eq:kernel}, evenness of $K_\alpha$, and
$C\preceq H_\alpha$ reduce the result to Theorem~\ref{thm:measure}.
\end{proof}

For finite paths and cycles, Lemma~\ref{lem:measures} gives
$\Lambda=4$, $K=1$, $\beta=1/2$, hence $A=2$ after the logarithmic
change of variables. Corollary~\ref{cor:cdf} proves
\eqref{eq:quantitative}. This proves the upper half of
Theorem~\ref{thm:main} for all paths and cycles, since those are exactly
the connected graphs of maximum degree at most two that have an edge.
In particular, arbitrary functions $n=n(\varepsilon)$ and
$p=p(\varepsilon)$ are covered by the same bound.

\section{Centered finite paths attain the coefficient}
We give the lower transfer explicitly, using convergence only in the
direction that supplies examples. Let $H_N$ be the relative derivative
for the central edge of the path on $2N$ vertices.
Only the odd indices contribute in Lemma~\ref{lem:measures}, each
with weight $1/N$. In angle coordinates their nodes are
\[
 \theta_j=\frac{(2j-1)\pi}{2N},\quad 1\le j\le N,
 \qquad \lambda(\theta)=4\sin^2(\theta/2).
\]
Embed their cyclic spaces as the step functions on $N$ equal cells
in $L^2((0,\pi),d\theta/\pi)$. The embedded derivative kernel is the
step approximation to
$K_\alpha(\log\lambda(\theta)-\log\lambda(\phi))$.
The kernels converge almost everywhere and are bounded by $\alpha$,
so dominated convergence gives Hilbert--Schmidt, hence operator-norm,
convergence to an operator $H_\infty$.
It acts on $L^2(\nu_\infty)$, where
\[
 d\nu_\infty(\lambda)
        =\frac{d\lambda}{\pi\sqrt{\lambda(4-\lambda)}},
                      \qquad 0<\lambda<4.
\]

On the low-energy compression $0<\lambda<1$, pass to
$x=\log\lambda<0$. Its kernel in Lebesgue measure becomes
\[
 \sqrt{r(x)r(y)}K_\alpha(x-y),\qquad
 r(x)=\frac{e^{x/2}}{\pi\sqrt{4-e^x}}
      =e^{x/2}v(x)^2,
\]
where $v$ is bounded and boundedly invertible on $(-\infty,0)$.
Equation~\eqref{eq:fourier} shows that
$\widehat K_\alpha/\widehat K_{1/2}$ is bounded above and below
by positive constants depending only on $\alpha$. Fourier quadratic
forms, compression, and multiplication by $e^{x/4}$ therefore give
a Loewner comparison of this compression with $M_v S M_v$, where
\[
                     S(x,y)=e^{(x+y)/4}K_{1/2}(x-y),\qquad x,y<0.
\]
Min--max and bounded invertibility of $M_v$ then imply
$\lambda_j(H_\infty)\ge c_\alpha\lambda_j(S)$ for some
$c_\alpha>0$. This multiplier comparison is a min--max argument,
not an entrywise kernel comparison.

Since $K_{1/2}(s)=1/(2\cosh(s/4))$, the unitary change
$u=e^{x/2}$ turns $S$ into the positive operator on $L^2(0,1)$
with kernel
\[
                              \frac{2\sqrt{uv}}{u+v}.
\]
Laptev's theorem~\cite[Theorem (3), p.~1039]{Laptev} gives count
$\log^2(1/\varepsilon)/(2\pi^2)$ for this kernel: in his notation
the denominator exponent is $a=1$, the numerator exponents are
$\beta=\gamma=1/2$, and
$\rho=\beta+\gamma-a+1=1$. All hypotheses hold; the factor two
only shifts the logarithmic threshold by a constant. Consequently
\begin{equation}\label{eq:inflower}
 \liminf_{\varepsilon\downarrow0}
 \frac{\N_{H_\infty}(\varepsilon)}{\log^2(1/\varepsilon)}
                          \ge\frac1{2\pi^2}.
\end{equation}

Each edited central edge is a bridge, so its exact correction satisfies
$C_N\succeq H_N/2$ by \eqref{eq:sandwich}.
For every fixed strict threshold, operator-norm convergence of the
embedded $H_N$ implies convergence of each ordered eigenvalue.
Thus, for every $\varepsilon>0$,
\[
 W_{\mathrm{path},\alpha}(\varepsilon)
 \ge\sup_N\N_{H_N}(2\varepsilon)
 \ge\N_{H_\infty}(2\varepsilon).
\]
Indeed each finite list strictly above $2\varepsilon$ in the limit
stays above it for all sufficiently large $N$.
Together with \eqref{eq:inflower} and the uniform upper bound this
proves Theorem~\ref{thm:main}.

\section{Scope and source checks}
The logarithmic localization and Fourier-decay method has classical
antecedents. Laptev~\cite[Sections 3--4]{Laptev} transforms weighted
Cauchy operators into weighted logarithmic convolutions, partitions
the half-line into cells, and controls diagonal and off-diagonal
contributions; his paper explicitly builds on Widom's method.
An earlier internal note,
\nolinkurl{matmul/uniform-alpha-channel-search.txt}, Sections~6--8,
already uses periodic cell kernels, positive Fourier decompositions
for arbitrary spectral measures, correction of the periodization
images, and optimization of cell masses for uniform Schatten bounds.
The internal \nolinkurl{channel-classification-priority-audit.txt}
expressly credits these Laptev--Widom antecedents.
Accordingly, no new periodization or phase-volume principle is claimed
here. These notes establish the explicit measure-tail count bound,
its stated uniformity and remainder, and its sharp finite-path/cycle
application by a self-contained version of that method.

The upper bound requires no path eigenvalue monotonicity, quadrature
ordering, or comparison of a finite matrix below its continuum limit.
The atomic measures are handled directly. Thus edges at fixed distance
from an endpoint, edges whose distances to the endpoints diverge at
arbitrary tolerance-dependent rates, and growing cycles all obey the
same upper coefficient. The fixed-endpoint limiting coefficient
$1/(6\pi^2)$ from the earlier session is compatible with this supremal
statement and is not used in this proof.

The uniform rank lemma and its block summation were independently
checked in three internal parallel reviews. One written check is
\texttt{periodization-independent-audit.tex} in this directory.
These are internal proof checks, not external peer review or formal
verification. No claim about historical priority is made.

For the lower bound, the primary Laptev PDF was inspected at
pp.~1038--1039, including the kernel, the parameter hypotheses, the
strict count convention, and Theorem~(3). Copies of these rendered
pages and the PDF are in \texttt{sources/}. For the Fourier symbol,
the DLMF table entry and its normalized Fourier convention were
inspected and rescaled to \eqref{eq:fourier}. The upper proof otherwise
uses only the explicit expansions and estimates given here.
All new artifacts from this session lie in the present directory;
the existing papers are unchanged.

\begin{thebibliography}{2}
\bibitem{Laptev}
A.~A.~Laptev, \emph{Spectral asymptotic behavior of a class of integral
operators}, Mathematical Notes 16 (1974), 1038--1043.
Theorem~(3), p.~1039.
\url{https://doi.org/10.1007/BF01149794}.
Author-hosted primary text:
\url{https://www.ma.ic.ac.uk/~alaptev/Papers/MatZam74.pdf}.
\bibitem{DLMF}
NIST Digital Library of Mathematical Functions, Table~1.14.1,
Fourier transforms, the $\sinh(at)/\sinh(\pi t)$ entry.
\url{https://dlmf.nist.gov/1.14.T1}.
\end{thebibliography}
\end{document}
